% Options for packages loaded elsewhere
\PassOptionsToPackage{unicode}{hyperref}
\PassOptionsToPackage{hyphens}{url}
\PassOptionsToPackage{dvipsnames,svgnames,x11names}{xcolor}
\documentclass[
  11pt,
]{article}
\usepackage{xcolor}
\usepackage[a4paper,margin=2.4cm]{geometry}
\usepackage{amsmath,amssymb}
\setcounter{secnumdepth}{-\maxdimen} % remove section numbering
\usepackage{iftex}
\ifPDFTeX
  \usepackage[T1]{fontenc}
  \usepackage[utf8]{inputenc}
  \usepackage{textcomp} % provide euro and other symbols
\else % if luatex or xetex
  \usepackage{unicode-math} % this also loads fontspec
  \defaultfontfeatures{Scale=MatchLowercase}
  \defaultfontfeatures[\rmfamily]{Ligatures=TeX,Scale=1}
\fi
\usepackage{lmodern}
\ifPDFTeX\else
  % xetex/luatex font selection
\fi
% Use upquote if available, for straight quotes in verbatim environments
\IfFileExists{upquote.sty}{\usepackage{upquote}}{}
\IfFileExists{microtype.sty}{% use microtype if available
  \usepackage[]{microtype}
  \UseMicrotypeSet[protrusion]{basicmath} % disable protrusion for tt fonts
}{}
\makeatletter
\@ifundefined{KOMAClassName}{% if non-KOMA class
  \IfFileExists{parskip.sty}{%
    \usepackage{parskip}
  }{% else
    \setlength{\parindent}{0pt}
    \setlength{\parskip}{6pt plus 2pt minus 1pt}}
}{% if KOMA class
  \KOMAoptions{parskip=half}}
\makeatother
\usepackage{longtable,booktabs,array}
\newcounter{none} % for unnumbered tables
\usepackage{calc} % for calculating minipage widths
% Correct order of tables after \paragraph or \subparagraph
\usepackage{etoolbox}
\makeatletter
\patchcmd\longtable{\par}{\if@noskipsec\mbox{}\fi\par}{}{}
\makeatother
% Allow footnotes in longtable head/foot
\IfFileExists{footnotehyper.sty}{\usepackage{footnotehyper}}{\usepackage{footnote}}
\makesavenoteenv{longtable}
\setlength{\emergencystretch}{3em} % prevent overfull lines
\providecommand{\tightlist}{%
  \setlength{\itemsep}{0pt}\setlength{\parskip}{0pt}}
% ─────────────────────────────────────────────────────────────────────────
% latex-header.tex — academic house style for the project's design docs.
% Injected into a standalone .tex by md2tex.sh (pandoc --include-in-header -s).
% Target: Overleaf with the default pdfLaTeX engine (no system fonts, no
% fontspec). Classic scholarly look: Palatino (newpx) body, navy headings,
% booktabs tables (pandoc emits these by default), microtype, A4.
% ─────────────────────────────────────────────────────────────────────────

% Fonts — Palatino-style serif via newpx (bundled with TeX Live / Overleaf).
\usepackage[T1]{fontenc}
\usepackage{newpxtext}
\usepackage{newpxmath}
\usepackage[scaled=0.92]{helvet}   % Helvetica-like sans
\usepackage{inconsolata}           % clean monospace

\usepackage{microtype}             % refined spacing / justification
\usepackage{xcolor}
\definecolor{headingnavy}{HTML}{1F3864}
\definecolor{rulegray}{HTML}{2F2F2F}

% Navy, serif section headings (matches the Word template).
\usepackage{titlesec}
\titleformat{\section}{\Large\bfseries\color{headingnavy}}{\thesection}{0.6em}{}
\titleformat{\subsection}{\large\bfseries\color{headingnavy}}{\thesubsection}{0.6em}{}
\titleformat{\subsubsection}{\normalsize\bfseries\color{headingnavy}}{\thesubsubsection}{0.6em}{}
\titlespacing*{\section}{0pt}{1.4ex plus 0.4ex}{0.7ex}
\titlespacing*{\subsection}{0pt}{1.1ex plus 0.3ex}{0.5ex}

% Booktabs rules in the same gray as the Word version.
\usepackage{booktabs}
\usepackage{colortbl}              % enables \arrayrulecolor
\arrayrulecolor{rulegray}
\setlength{\heavyrulewidth}{1.1pt}
\setlength{\lightrulewidth}{0.5pt}
\renewcommand{\arraystretch}{1.25}
\usepackage{ragged2e}

% Unicode glyph that pdfLaTeX/newpx doesn't render on its own.
% (≤ and ≥ are already provided, so we don't redefine them here.)
\usepackage{newunicodechar}
\newunicodechar{→}{\ensuremath{\rightarrow}}

% NOTE: link colours are set via pandoc -V variables (colorlinks/linkcolor/
% urlcolor) — NOT \hypersetup here, because pandoc loads hyperref *after*
% this header, so calling \hypersetup here would be an undefined command.

% Report-style spacing: blank line between paragraphs, no indent.
\usepackage{parskip}
\usepackage{bookmark}
\IfFileExists{xurl.sty}{\usepackage{xurl}}{} % add URL line breaks if available
\urlstyle{same}
\hypersetup{
  colorlinks=true,
  linkcolor={headingnavy},
  filecolor={Maroon},
  citecolor={Blue},
  urlcolor={headingnavy},
  pdfcreator={LaTeX via pandoc}}

\author{}
\date{}

\begin{document}

\section{User Stories}\label{user-stories}

\textbf{Project:} RAG-Based Retrieval for Archival Systems\\
\textbf{Version:} 1.0\\
\textbf{Date:} 2026-06-01\\
\textbf{Status:} Draft

\subsection{Roles}\label{roles}

{\def\LTcaptype{none} % do not increment counter
\begin{longtable}[]{@{}
  >{\raggedright\arraybackslash}p{(\linewidth - 2\tabcolsep) * \real{0.5000}}
  >{\raggedright\arraybackslash}p{(\linewidth - 2\tabcolsep) * \real{0.5000}}@{}}
\toprule\noalign{}
\begin{minipage}[b]{\linewidth}\raggedright
Role
\end{minipage} & \begin{minipage}[b]{\linewidth}\raggedright
Description
\end{minipage} \\
\midrule\noalign{}
\endhead
\bottomrule\noalign{}
\endlastfoot
\textbf{Researcher} & Academic or graduate student querying the archive
for primary sources, historical data, or literature. Needs precise,
citable answers. \\
\textbf{Professional} & Practitioner (e.g., lawyer, journalist, museum
curator) who needs to quickly locate specific archival records or
policies. Needs fast, trustworthy responses. \\
\textbf{General Public User} & Member of the public exploring archival
content. Needs accessible, plain-language responses without domain
expertise. \\
\textbf{Archivist / Admin} & Institutional staff managing the Islandora
repository, approving ingestions, and maintaining collection quality. \\
\textbf{System Administrator} & Engineer responsible for deploying,
monitoring, and maintaining the RAG infrastructure. \\
\end{longtable}
}

\subsection{Epic 1: Natural Language Document
Retrieval}\label{epic-1-natural-language-document-retrieval}

\emph{Users can ask questions in natural language and receive grounded,
cited answers from the archival collection.}

\subsubsection{US-01: Ask a Natural Language
Question}\label{us-01-ask-a-natural-language-question}

\textbf{As a} researcher,\\
\textbf{I want to} type a question in natural language\\
\textbf{so that} I can find relevant archival information without
learning a query syntax or search operators.

\textbf{Priority:} High\\
\textbf{Effort:} L

\textbf{Acceptance Criteria:}

\begin{itemize}
\tightlist
\item
  The chat interface accepts free-text input
\item
  A response is returned within 10 seconds
\item
  The response directly addresses the question's intent
\end{itemize}

\subsubsection{US-02: Receive Source-Attributed
Response}\label{us-02-receive-source-attributed-response}

\textbf{As a} researcher,\\
\textbf{I want} every response to include citations naming the source
document, section, and page\\
\textbf{so that} I can verify the answer independently and use it in
academic work.

\textbf{Priority:} High\\
\textbf{Effort:} M

\textbf{Acceptance Criteria:}

\begin{itemize}
\tightlist
\item
  Each factual claim in the response is accompanied by a citation
\item
  The citation identifies at minimum the document name and section/page
\item
  Clicking or tapping a citation opens or highlights the source passage
\end{itemize}

\subsubsection{US-03: Ask Follow-Up
Questions}\label{us-03-ask-follow-up-questions}

\textbf{As a} researcher,\\
\textbf{I want to} ask follow-up questions that reference my earlier
query\\
\textbf{so that} I can explore a topic iteratively without restating
context every time.

\textbf{Priority:} Medium\\
\textbf{Effort:} M

\textbf{Acceptance Criteria:}

\begin{itemize}
\tightlist
\item
  A follow-up such as ``Can you elaborate on the second point?''
  resolves correctly to prior context
\item
  Session history is maintained for the duration of the session
\item
  The system does not mix context from different sessions
\end{itemize}

\subsubsection{US-04: Query in My Own
Language}\label{us-04-query-in-my-own-language}

\textbf{As a} general public user whose primary language is not
English,\\
\textbf{I want to} submit queries in my native language (e.g., German)\\
\textbf{so that} I can access archival content without being fluent in
the archive's primary language.

\textbf{Priority:} Medium\\
\textbf{Effort:} L

\textbf{Acceptance Criteria:}

\begin{itemize}
\tightlist
\item
  Queries in supported languages are correctly understood and answered
\item
  Results from source documents in different languages are returned when
  relevant
\item
  The response language matches the query language
\end{itemize}

\subsubsection{US-05: Filter by Document
Metadata}\label{us-05-filter-by-document-metadata}

\textbf{As a} professional,\\
\textbf{I want to} narrow my search by date range, collection, or
document type\\
\textbf{so that} I get results from the specific subset of the archive
relevant to my task.

\textbf{Priority:} Medium\\
\textbf{Effort:} M

\textbf{Acceptance Criteria:}

\begin{itemize}
\tightlist
\item
  The interface exposes date range, collection tag, and document type
  filters
\item
  Retrieval respects applied filters without the user embedding them in
  query text
\item
  Filters can be combined
\end{itemize}

\subsubsection{US-06: Get Results Despite Historical
Terminology}\label{us-06-get-results-despite-historical-terminology}

\textbf{As a} researcher working with pre-20th century archives,\\
\textbf{I want} the system to find relevant documents even when I use
modern vocabulary\\
\textbf{so that} vocabulary mismatch between my query and historical
source language does not cause gaps in results.

\textbf{Priority:} High\\
\textbf{Effort:} L

\textbf{Acceptance Criteria:}

\begin{itemize}
\tightlist
\item
  A query for ``aviation'' returns documents containing ``flying
  machines'' where relevant
\item
  Hybrid retrieval (dense + sparse) is active by default
\item
  No user configuration is required to benefit from semantic bridging
\end{itemize}

\subsubsection{US-07: Query from the Drupal
Front-End}\label{us-07-query-from-the-drupal-front-end}

\textbf{As a} general public user,\\
\textbf{I want to} query the archive directly within the Drupal
repository site\\
\textbf{so that} I can access institutional knowledge without leaving
the platform I'm already browsing.

\textbf{Priority:} Medium\\
\textbf{Effort:} L

\textbf{Acceptance Criteria:}

\begin{itemize}
\tightlist
\item
  The Drupal front-end provides a chat interface for submitting queries
\item
  Responses and citations render correctly within the Drupal page
\item
  The front-end connects to the backend query flow orchestrated by
  Node-RED
\end{itemize}

\subsection{Epic 2: Document Ingestion \& Collection
Management}\label{epic-2-document-ingestion-collection-management}

\emph{Archivists can add new documents to the collection and have them
become immediately queryable.}

\subsubsection{US-08: Ingest a New Document via
Islandora}\label{us-08-ingest-a-new-document-via-islandora}

\textbf{As an} archivist,\\
\textbf{I want to} add a new digitised document to Islandora\\
\textbf{so that} it becomes searchable by all users without requiring
developer intervention.

\textbf{Priority:} High\\
\textbf{Effort:} L

\textbf{Acceptance Criteria:}

\begin{itemize}
\tightlist
\item
  Uploading a document to Islandora triggers automatic pipeline
  processing
\item
  The document's text is extracted, chunked, embedded, and indexed in
  ChromaDB
\item
  The document is queryable within 15 minutes for a 100-page document
\end{itemize}

\subsubsection{US-09: Preserve Document Structure During
Ingestion}\label{us-09-preserve-document-structure-during-ingestion}

\textbf{As an} archivist,\\
\textbf{I want} document section headings, hierarchy, and metadata
preserved in each chunk\\
\textbf{so that} retrieved passages can be correctly understood in
context by the LLM.

\textbf{Priority:} High\\
\textbf{Effort:} M

\textbf{Acceptance Criteria:}

\begin{itemize}
\tightlist
\item
  Each stored chunk includes \texttt{context\_heading},
  \texttt{context\_content}, \texttt{context\_links}, and metadata
  fields
\item
  No chunk boundary splits a sentence mid-phrase or separates a table
  row from its header
\item
  Structure-aware chunking is used, not naive character-count splitting
\end{itemize}

\subsubsection{US-10: Track Ingestion
Status}\label{us-10-track-ingestion-status}

\textbf{As an} archivist,\\
\textbf{I want to} see the ingestion status of each document I submit\\
\textbf{so that} I know when it is available for queries and can
investigate failures.

\textbf{Priority:} Low\\
\textbf{Effort:} M

\textbf{Acceptance Criteria:}

\begin{itemize}
\tightlist
\item
  The admin interface shows per-document status: pending, processing,
  indexed, failed
\item
  Failed ingestions display a human-readable error message
\item
  Successful ingestions show a timestamp and chunk count
\end{itemize}

\subsubsection{US-11: Remove a Document from the
Index}\label{us-11-remove-a-document-from-the-index}

\textbf{As an} archivist,\\
\textbf{I want to} remove a document from the retrieval index\\
\textbf{so that} users cannot receive responses grounded in outdated,
incorrect, or restricted content.

\textbf{Priority:} Medium\\
\textbf{Effort:} M

\textbf{Acceptance Criteria:}

\begin{itemize}
\tightlist
\item
  Removing a document triggers deletion of its chunks from ChromaDB
\item
  Subsequent queries no longer return content from that document
\item
  The operation does not require a full re-index of the collection
\end{itemize}

\subsection{Epic 3: System Operations}\label{epic-3-system-operations}

\emph{The system administrator can deploy and maintain the pipeline on a
single GPU node.}

\subsubsection{US-12: Deploy on a Single 24 GB
GPU}\label{us-12-deploy-on-a-single-24-gb-gpu}

\textbf{As a} system administrator,\\
\textbf{I want to} run the full stack on a single NVIDIA L4 or RTX
4090\\
\textbf{so that} the prototype can be deployed without a multi-GPU
cluster or cloud infrastructure.

\textbf{Priority:} High\\
\textbf{Effort:} L

\textbf{Acceptance Criteria:}

\begin{itemize}
\tightlist
\item
  LLM inference, embedding model, and reranker run within 24 GB VRAM
\item
  VRAM usage is documented for the recommended model sizes (7B--14B
  quantized)
\item
  RAGDoll-style offloading is configured to handle memory contention
\end{itemize}

\subsubsection{US-13: Keep All Data
On-Premises}\label{us-13-keep-all-data-on-premises}

\textbf{As a} system administrator,\\
\textbf{I want} all inference and document storage to remain
on-premises\\
\textbf{so that} institutional data sovereignty and privacy requirements
are met.

\textbf{Priority:} High\\
\textbf{Effort:} M

\textbf{Acceptance Criteria:}

\begin{itemize}
\tightlist
\item
  No document content, query text, or generated response is transmitted
  to an external API
\item
  LLM runs via OLLAMA or vLLM on local hardware
\item
  All dependencies are open-source with no mandatory cloud licensing
\end{itemize}

\end{document}
